\section{Problem Identification}\label{sec:identification}

An investor has one unit of capital and must divide it among the stocks of the VN100 index, the
hundred largest and most actively traded companies listed in Vietnam. Ninety-eight of those
tickers survive the data cleaning of \cref{sec:data}, and the decision is informed by four
years of daily closing prices. What fraction of capital should go into each stock? This is the
general problem of \textbf{portfolio optimization}: choosing capital weights that are, in some
precise sense, the best response to a set of assets' estimated returns and risks.

\textbf{There is no single formula that is best for every market.} The literature offers several
competing answers to what ``best'' should mean --- plain mean--variance efficiency
\citep{markowitz1952}, robustness to estimation error \citep{goldfarb2003}, an explicit penalty
on the number of positions held \citep{tibshirani1996, brodie2009} --- and which of these actually
helps depends on the statistical character of the specific universe being invested in: how noisy
its mean returns are, how correlated its constituents are, how expensive its transaction costs
are. A formula tuned on, say, US large-cap equities carries no guarantee of transferring to a
100-stock Vietnamese universe with a different correlation structure and a different cost profile.
The approach taken here is accordingly not to commit upfront to one fixed formula and assume it is
right, but to state \textbf{one general objective} whose competing terms can each be switched on,
off, or tuned --- collapsing to several named strategies as special cases --- and to let the data
decide, empirically, which configuration of that objective actually performs best on \emph{this}
universe.

The textbook starting point for that general objective is Markowitz mean--variance analysis
\citep{markowitz1952}: estimate each stock's average return and the covariances between stocks,
then trade expected return against portfolio variance. That answer alone is incomplete for three
reasons, and each motivates a term the general formula must be able to switch on:

\textbf{The inputs are estimates, and the mean is a bad one.} An average return computed from a
few years of daily data is noisy. An optimizer handed such an estimate will concentrate capital
in whichever stocks happened to drift upward during the sample, because it cannot distinguish a
genuine tendency from an accident of the period; the resulting portfolio is optimal for the past
and arbitrary for the future. This is the standard explanation for why optimized portfolios so
often fail to beat a naive equal-weight benchmark out of sample \citep{demiguel2009}. A useful
model should therefore be able to ask not only what is best if the estimate is right, but what
remains acceptable if it is wrong.

\textbf{Holding ninety-eight stocks is expensive.} A mathematically optimal portfolio typically
assigns a nonzero weight to every available asset, many of them a fraction of a percent. Each
position must be bought, rebalanced, monitored, and sold, and each transaction costs money. An
investor generally prefers ten or twenty names to ninety-eight even at some cost in theoretical
optimality --- a preference worth encoding as an optional term rather than ignoring or forcing on
unconditionally.

\textbf{Risk still matters.} Variance of the realized portfolio return remains a genuine concern,
and the classical quadratic penalty remains the natural way to express it, in every
configuration.

The goal is therefore twofold: first, to state one model, balancing up to four competing
objectives at once --- expected return, variance, sensitivity to a mis-estimated mean, and the
number of positions held --- general enough that fixing or dropping its terms reproduces several
named strategies rather than requiring a separate formula for each; second, to treat the choice
\emph{among} those configurations as an empirical question specific to VN100 rather than a
theoretical one settled in advance, answered by walk-forward testing several configurations
against each other on this exact universe and reporting which one actually wins here.
\Cref{sec:modeling} states and classifies the general model, \cref{sec:correctness} argues it
represents the problem faithfully and is convex and second-order-cone representable,
\cref{sec:data} describes the data, \cref{sec:insample} solves the model directly with
CVXPY/CLARABEL and reports in-sample results, \cref{sec:backtest} turns the choice of
configuration into that empirical test --- six named variants, walk-forward out of sample, on the
VN100 universe specifically --- and \cref{sec:conclusion} discusses efficiency and limitations.
